\documentclass{article}
\usepackage{amsmath,amssymb}
\usepackage{xurl}
\usepackage[hidelinks]{hyperref}
% Shared commands for notes in docs/paper-gaps/.

\DeclareUnicodeCharacter{2080}{\ifmmode{}_{0}\else\textsubscript{0}\fi}
\DeclareUnicodeCharacter{2081}{\ifmmode{}_{1}\else\textsubscript{1}\fi}
\DeclareUnicodeCharacter{2082}{\ifmmode{}_{2}\else\textsubscript{2}\fi}
\DeclareUnicodeCharacter{2099}{\ifmmode{}_{n}\else\textsubscript{n}\fi}

\newcommand{\Lean}{\textsc{Lean}}
\ExplSyntaxOn
\NewDocumentCommand{\leanid}{m}{
  \ifmmode
    \textsf{\detokenize{#1}}
  \else
    \group_begin:
    \urlstyle{sf}
    \tl_set:Nn \l_tmpa_tl {#1}
    \tl_replace_all:Nnn \l_tmpa_tl {\_} {_}
    \exp_args:NV \path \l_tmpa_tl
    \group_end:
  \fi
}
\ExplSyntaxOff
\providecommand{\tr}{\operatorname{tr}}
\newcommand{\tnleanIssueBase}{https://github.com/LionSR/TNLean/issues/}
\newcommand{\tnleanPrBase}{https://github.com/LionSR/TNLean/pull/}
\newcommand{\tnleanDocBase}{https://sirui-lu.com/TNLean/docs/}
\newcommand{\ghissue}[1]{\href{\tnleanIssueBase#1}{GitHub issue \##1}}
\newcommand{\ghpr}[1]{\href{\tnleanPrBase#1}{PR \##1}}
\newcommand{\doclink}[1]{\href{\tnleanDocBase#1.html}{\path{#1}}}

% Dirac notation and the identity operator, as in blueprint/src/macros/common.tex.
% \providecommand keeps note-local definitions authoritative.
\usepackage{dsfont}
\providecommand{\ket}[1]{|#1\rangle}
\providecommand{\bra}[1]{\langle #1|}
\providecommand{\braket}[2]{\langle #1 | #2 \rangle}
\providecommand{\ketbra}[2]{|#1\rangle\!\langle #2|}
\providecommand{\Id}{\mathds{1}}
\providecommand{\one}{\mathds{1}}
\providecommand{\C}{\mathbb{C}}
\providecommand{\MN}[1]{M_{#1}(\C)}
\providecommand{\MD}{M_D(\C)}

% Verdict marker: \gapnote{<kind>}{<status>}. Kind is the policy
% classification (clarification, local-correction, scope-restriction,
% unfaithful, false-source, open-gap); status is open, wip (elimination
% underway), resolved, or historical. Machine-read by the site index;
% prints a verdict line.
\newcommand{\gapnote}[2]{\par\noindent\textbf{Verdict:} #1 (#2).\par}

\providecommand{\leanid}[1]{\textsf{\detokenize{#1}}}
\providecommand{\gapnote}[2]{\par\noindent\textbf{Verdict:} #1 (#2).\par}
\title{Numerical Power Decay of the Selected Charge Energy}
\author{}
\date{2026-10-14}
\begin{document}
\maketitle
\gapnote{scope-restriction}{open}

\section{Precise numerical conclusion}

Fix the final scanner exponents
$\omega=1/20000$, $\ell=1/100000$, $\kappa=1/1000000$
and $\nu=1/1000$. For an integer $n\geq1$, retain the literal
separation ceiling $D=\lceil n^\kappa\rceil$ and define
\begin{align}
    E_n(W,x,p)
      &=CW^2n^\ell D^2(\log n)^{C_\ell}
          \left((1-p)x^{1/8}+(1-p)n^{-P}+p\right).
    \label{eq:selected_energy_scope_cap}
\end{align}
Here $x$ is the normalized entropy defect. Fix $C\geq0$,
$C_\ell\in\mathbb R$, $g>0$ and $P\geq\nu$.
There is a threshold $N\geq2$, selected after these fixed parameters
and before all scalar data, such that for every integer $n\geq N$,
\begin{align}
    1\leq W\leq n^\omega,\qquad
    0\leq x\leq n^{-8\nu},\qquad
    n^{-\nu}\leq p\leq2n^{-\nu},\qquad
    e\leq n^{-1000}
    \label{eq:selected_energy_scope_hypotheses}
\end{align}
imply
\begin{align}
    \frac{e+E_n(W,x,p)}{g/2}&\leq n^{-1/2000},&
    \frac{e+E_n(W,x,p)+r}{g/2}
      &\leq n^{-1/2000}+\frac r{g/2}
    \label{eq:selected_energy_scope_conclusion}
\end{align}
for every real $r$. The source's stronger restrictions $C\geq1$,
$C_\ell\geq0$ and $0\leq e\leq\min(n^{-1000},g/4)$ are therefore
covered. No assumption on $r$ is needed for the displayed inequality.
\footnote{The declarations are
  \leanid{TNLean.PEPS.AreaLaw.Scan.exists_final_selectedEnergy_le_rpow}
  and
  \leanid{TNLean.PEPS.AreaLaw.Scan.exists_final_selectedEnergy_add_remainder_le_rpow}.
  Their eventual form is
  \leanid{TNLean.PEPS.AreaLaw.Scan.eventually_final_selectedEnergy_le_rpow}.}

\section{Power arithmetic and sufficient tail condition}

The matched density rate is $8\nu=1/125$. Its eighth root is at most
$n^{-\nu}$. The explicit sufficient condition $P\geq\nu=1/1000$
puts the bad-old-history term at the same rate. Together with the full
interpolation interval in (\ref{eq:selected_energy_scope_hypotheses}),
these give a bracket at most $4n^{-\nu}$.
The ceiling contributes $D^2\leq4n^{2\kappa}$, so
\begin{align}
    E_n(W,x,p)&\leq16C n^{2\omega+\ell+2\kappa-\nu}(\log n)^{C_\ell},&
    2\omega+\ell+2\kappa-\nu&=-\frac{111}{125000}.
    \label{eq:selected_energy_scope_power}
\end{align}
The exact strict margin to $-1/2000$ is $97/250000>0$.
Every fixed logarithmic power is absorbed by this margin, while
$n^{-1000}$ is absorbed separately. Each contribution is bounded by
$(g/4)n^{-1/2000}$, so their sum divided by $g/2$ satisfies
(\ref{eq:selected_energy_scope_conclusion}).
\footnote{The pointwise bound and exact arithmetic are
  \leanid{TNLean.PEPS.AreaLaw.Scan.finalSelectedEnergyCap_le}
  and \leanid{TNLean.PEPS.AreaLaw.Scan.final_selectedEnergy_margins}.}

The condition $P\geq\nu$ is sufficient, rather than asserted to be
minimal. The source permits an arbitrary fixed inverse-power tail in the
old-history estimate; this specialization uses that freedom explicitly.
The energy expression and final parameter values are from
\cite{OpenAI2026AreaLaw}, \path{08-scanner.tex}, lines 492--511 and
839--865, revision
\texttt{adc7f1241b42e322a6451854ab7e4b4c146bf78a}.

\section{Remaining physical obligations}

This result proves the numerical implication
(\ref{eq:selected_energy_scope_hypotheses})--
(\ref{eq:selected_energy_scope_conclusion}). It does not derive
$x\leq n^{-8\nu}$ from physical entropy data, construct a selected
charge point, or show that the physical energy defect is bounded by
$(e+E_n+r)/(g/2)$. To obtain a physical conclusion, the stronger density
estimate and the energy comparison must hold for the same scanner,
replica labels, selected round and interpolation parameter.
Those same-point obligations are not premises hidden inside the scalar
proof: they are separate inputs required for its application.

If a physical remainder $r_k$ tends to zero with the replica count,
its contribution remains $r_k/(g/2)$ in
(\ref{eq:selected_energy_scope_conclusion}). This result neither absorbs
that remainder into a system-size power nor supplies a uniform rate in
both variables. The same-point sharp comparison, the remaining entropy
bootstrap, geometric estimates and limit and continuity arguments are
not completed here. No unrestricted scanner proposition or area law is
asserted. The scope restriction remains open.

\bibliographystyle{alpha}
\bibliography{references}
\end{document}
